% SPDX-License-Identifier: Apache-2.0 OR MIT
\documentclass[conference]{IEEEtran}
\IEEEoverridecommandlockouts

% T1 keeps the legacy Times metrics IEEEtran is designed around; without it a
% Unicode engine falls back to substituted shapes and sets the tables wider.
\usepackage[T1]{fontenc}
\usepackage{cite}
\usepackage{amsmath}
\usepackage{amssymb}
\usepackage{booktabs}
\usepackage{graphicx}
\usepackage{url}
\usepackage[hidelinks]{hyperref}

\def\BibTeX{{\rm B\kern-.05em{\sc i\kern-.025em b}\kern-.08em
    T\kern-.1667em\lower.7ex\hbox{E}\kern-.125emX}}

% Long C++ identifiers are unbreakable boxes in a 3.5-inch column. Giving TeX
% more tolerance lets it choose a break sequence that avoids one very loose
% line rather than distributing the slack.
\tolerance=2000
\emergencystretch=2em

\begin{document}

\title{Aparajita: Branchless SIMD Search and\\
       Append-Only Nodes for LSM-Tree MemTables}

\author{\IEEEauthorblockN{Parth Kumar Sinha}
\IEEEauthorblockA{Independent Researcher\\
Oxford, United Kingdom\\
ORCID: 0009-0002-3120-9301}}

\maketitle

\begin{abstract}
The in-memory write buffer of an LSM-tree key-value store is searched on every
read and written on every insert, so its cost per operation sets a floor on the
whole engine. RocksDB uses a concurrent skip list, which chases pointers across
independently allocated nodes and evaluates one data-dependent branch per key
comparison. We present Aparajita, a MemTable representation that replaces the
skip list with a list of cache-line-sized nodes, each holding fifteen 32-bit
order-preserving key surrogates and a sentinel in one line, searched by a
branchless SIMD kernel. Three
design decisions carry the result. A relational vector compare over a sorted node
yields a mask whose population count is the lower bound directly, so ordered
search costs one compare, one movemask and one popcount with no branch. Surrogates
are taken after the node's shared prefix rather than from the start of the key,
without which an absolute surrogate takes one value across all 200{,}000 keys in
five of eight realistic distributions, including the keyspace this paper's own
evaluation runs on. Nodes are append-only, and the sorted order
over their slots is a 64-bit word, so an insert is two stores into a free slot
followed by one release store that publishes them.

We implement Aparajita as a RocksDB plugin selectable by name without patching
RocksDB sources, and evaluate it against the default skip list and VectorRep on a
12-core Emerald Rapids host. Point lookups over a resident MemTable are 22\% to
39\% faster than the skip list at 1, 4, 16 and 64 threads, non-overlapping across
five runs per configuration at every point but one, backed by 55\% fewer retired
instructions and 40\% fewer L1 misses per lookup. Ordered seeks are 15\% to 29\%
faster, but a seek followed by ten iterator steps is 3\% to 4\% slower, and the
representation charges 1.4 times the skip list's arena per key. The skip list is
48\% to 52\% slower on insert at the representation level, and Aparajita is 18.8\%
faster in single-threaded \texttt{db\_bench},
but the multi-threaded \texttt{db\_bench} write path is bounded by RocksDB's write
group rather than by the MemTable: an insert there retires over 22{,}000
instructions in both representations. We report that ceiling rather than a write
scaling claim the data does not support.
\end{abstract}

\begin{IEEEkeywords}
LSM-tree, key-value store, MemTable, SIMD, branch prediction, cache-conscious
data structures, RocksDB
\end{IEEEkeywords}

\section{Introduction}

An LSM-tree \cite{oneil1996lsm,luo2020survey} absorbs writes into an in-memory
buffer and flushes it to sorted files on disk when it fills. In RocksDB
\cite{dong2021rocksdb,dong2017space,rocksdb} and in LevelDB \cite{leveldb},
from which it descends, that buffer is the MemTable, and every read
consults it before any file, so it sits on both the write path and the read path
of the entire engine. Its default implementation is a concurrent skip list
\cite{pugh1990skip}.

A skip list is a good fit for the interface and a poor fit for the hardware. Each
node is separately allocated, so a search follows a chain of pointers whose
targets are unrelated in memory, and each hop costs a cache miss that the
prefetcher cannot anticipate because the address is not known until the previous
load returns. Each hop also compares one key against one bound, which is a
data-dependent branch, and the branch is unpredictable by construction: the whole
point of the comparison is that its outcome is not known in advance. Modern cores
are built to hide latency through parallelism, and a dependent chain of misses
followed by mispredicted branches defeats both mechanisms at once.

The alternative this paper explores is to make the unit of search a single cache
line. If sixteen 32-bit lanes fit in one 64-byte line and can be compared in one or
two vector instructions, then a search over the keys they hold costs one cache fill,
no pointer chase, and no branch at all. Reaching the right line still requires a
descent, so the design also has to make the descent cheap; a fast leaf reached by
an expensive walk is not an improvement.

Getting from that idea to a working RocksDB representation raises three problems
that shape the rest of the paper.

RocksDB keys are variable-length internal keys, not 32-bit integers, so a vector
lane holds a surrogate rather than a key. The surrogate has to preserve order, or
the same kernel cannot serve both point lookups and the ordered iteration that
RocksDB requires. It also has to discriminate, which turns out to rule out the
obvious choice of leading key bytes.

An ordered structure with fixed-capacity nodes needs an insert that keeps the node
sorted, and shifting keys inside a node is not a single atomic operation. A design
that publishes each insert through one release store therefore has to separate the
key material from the order over it.

The measurement itself is a problem. A MemTable representation is measured through
a benchmark harness that also exercises a write-ahead log, a block cache, a
version set and a write-thread group, and several of those cost more per operation
than the structure under test. Part of this paper's contribution is identifying
where the harness, rather than the representation, is what a number describes.

\subsection{Contributions}

\begin{itemize}
\item A cache-line node layout with a branchless ordered SIMD search kernel whose
      population count is the lower bound directly, measured at 4.03 cycles per
      probe on AVX-512 and 5.10 on AVX2 against 40.14 for the branchy scalar lower
      bound it replaces, with zero branch mispredictions per probe.
\item A prefix-relative order-preserving surrogate, motivated by a measurement
      showing that an absolute leading-bytes surrogate takes a single value across
      all 200{,}000 keys in five of eight realistic distributions, the keyspace
      the evaluation itself uses among them.
\item An append-only node whose sorted order is a 64-bit word of packed slot
      indices, making an insert one release store and reducing arena consumption
      from 445.4 to 100.6 bytes per key.
\item An eight-byte descent hint in the node header that replaces a virtual
      comparator call per tower hop with an integer compare in a line already
      loaded.
\item An evaluation against RocksDB's skip list and VectorRep covering point
      lookups, ordered seeks, inserts, arena cost and flush cost, that separates
      the representation's cost from the harness's, including two measurement
      errors we made and corrected and an exit criterion we could not meet.
\end{itemize}

\section{Background and related work}

\subsection{MemTable representations in RocksDB}
\label{sec:reps}

RocksDB exposes the MemTable through the \texttt{MemTableRep} interface and ships
several implementations. The default \texttt{SkipListRep} wraps a concurrent skip
list supporting lock-free reads and concurrent inserts. \texttt{VectorRep} appends
to a \texttt{std::vector} and sorts when the MemTable becomes immutable, which
makes its inserts very cheap and its reads on a mutable buffer very expensive,
because a point lookup on an unsorted vector has to sort a copy.
\texttt{HashSkipListRep} and \texttt{HashLinkListRep} bucket by prefix and are
useful only when the workload's access pattern matches the prefix extractor.

A reviewer familiar with RocksDB will ask why Aparajita is not VectorRep with
vector instructions. The answer is that VectorRep gives up ordering, and ordering
is what a MemTable iterator has to provide. Section~\ref{sec:vectorrep} measures
what that costs VectorRep in practice, and Section~\ref{sec:seeks} measures what
providing it costs us.

\subsection{Cache-conscious and SIMD search}

Making search structures fit the memory hierarchy is long-established. Rao and
Ross showed that removing pointers from B$^+$-tree nodes improves main-memory
search \cite{rao2000cache}. FAST laid out a tree so that its levels match cache
line, page and SIMD widths \cite{kim2010fast}. Schlegel et al.\ generalised binary
search to $k$-ary search driven by vector comparisons \cite{schlegel2009kary}, and
Zeuch et al.\ applied similar transformations to tree traversal
\cite{zeuch2014simd}. Masstree combines a trie of B$^+$-trees with careful
attention to cache line behaviour \cite{mao2012masstree}, the Adaptive Radix Tree
adapts node representations to density \cite{leis2013art}, and the Bw-Tree reaches
latch freedom through a mapping table and delta records \cite{levandoski2013bwtree}.

Masstree deserves separating out, because its keyslice is the closest prior art to
two of the mechanisms here. Masstree cuts a variable-length key into fixed
eight-byte slices and compares a slice as an integer, falling back to the full key
when slices tie, which is structurally what both our surrogate and our descent hint
do. The difference is where the fixed window sits. A Masstree slice is taken at a
depth fixed by the trie level a node occupies, so the bytes it covers are the same
for every node at that depth. Our surrogate is taken at an offset each node
computes for itself, the length of the prefix its own fifteen keys share, which is
what keeps a lane discriminating on a keyspace whose shared bytes run deeper than
a fixed window would skip. Section~\ref{sec:surrogate} measures the difference that
makes. The descent hint is the fixed-window case and takes the first eight bytes
outright, because it is compared across nodes rather than within one.

Aparajita differs from the rest of this line of work in what it is embedded in
rather than in the kernel itself. These structures are read-optimised indexes built once or
updated infrequently. A MemTable is written continuously by many threads while
being read, is destroyed after a few seconds of life, and must allocate from a
bump allocator that cannot free or reallocate. The kernel is the easy part; the
concurrency protocol and the allocation discipline around it are where the design
effort goes.

\subsection{Key compression and surrogates}

Compressing keys into fixed-width order-preserving codes is the subject of HOPE
\cite{zhang2020hope}, and hybrid indexes \cite{zhang2016hybrid} use compact
representations for cold data. Our surrogate is a much simpler object, four bytes
taken at a node-local offset, and it is a filter rather than a decision: every
match it produces is confirmed against RocksDB's comparator before it is used.
Section~\ref{sec:surrogate} explains why removing that confirmation would lose
keys rather than merely slow lookups down.

\section{Design}

\subsection{Node layout}

A node holds sixteen 32-bit lanes in exactly one 64-byte cache line, aligned so it
never straddles a line boundary. Fifteen of those lanes hold key surrogates and the
sixteenth is a permanent sentinel, for the reason Section~\ref{sec:append} gives;
the node's capacity is fifteen keys and we write it that way throughout. AVX2
covers the line in two 256-bit compares and AVX-512 in one, so a probe costs one
line fill and a fixed number of instructions independent of where the answer lies.
Figure~\ref{fig:layout} shows the two lines.

Everything else the node needs, the order word, the shared prefix length, the
level-0 successor pointer and the key pointers, begins at offset 64.
The separation is deliberate: a search that misses touches only the first line,
and the cold fields are reached only after the vector step has narrowed the answer
to a single candidate. The layout is enforced by static assertions rather than by
comment, because the search kernel reinterprets the surrogate array as a vector
and would silently misbehave if a field migrated into the first line.

\begin{figure}[t]
\centering\footnotesize
\setlength{\tabcolsep}{3pt}
\begin{tabular}{@{}rll@{}}
\toprule
byte & field & touched by \\
\midrule
\multicolumn{3}{@{}l}{\emph{\texttt{NodeData}, line 0: all a probe that misses reads}}\\
0--59   & surrogates $s_0 \dots s_{14}$          & vector compare \\
60--63  & sentinel lane, \texttt{0xFFFFFFFF}     & vector compare \\
\midrule
\multicolumn{3}{@{}l}{\emph{\texttt{NodeData}, line 1 on: read after the vector step}}\\
64--71  & order word, rank $i$ in nibble $i$     & every reader (acquire) \\
72--75  & \texttt{prefix\_len}                   & surrogate computation \\
80--87  & level-0 successor                      & descent, iteration \\
88--    & 16 key slots, 15 ever live             & comparator confirmation \\
\midrule
\multicolumn{3}{@{}l}{\emph{\texttt{ListNode}, its own line: the index above level 0}}\\
0--7    & payload pointer                        & swapped by a split \\
8--15   & \texttt{first\_hint}, 8 key bytes      & every tower hop \\
16--    & write lock, height, tower              & writers, descent \\
\bottomrule
\end{tabular}
\caption{Node layout. The split at offset 64 is what makes a probe that finds
nothing cost one line fill: the order word, the prefix length, the successor and
the key pointers are all reached only once the vector step has narrowed the answer
to a candidate. The successor and the order word live inside \texttt{NodeData}
rather than beside the tower so that a split can replace contents and successor in
one release store. Offsets are enforced by static assertion, since the search
kernel reinterprets bytes 0--63 as a vector.}
\label{fig:layout}
\end{figure}

\subsection{Branchless ordered search}
\label{sec:kernel}

The kernel is where the branch mispredictions go away, and the reason it works for
an ordered structure is a property of the mask rather than of the compare.

Let a node hold surrogates $s_0 \le s_1 \le \dots \le s_{15}$ and let $t$ be the
target. A relational compare produces the bit mask $m_i = [s_i < t]$. Because the
node is sorted, the set bits of $m$ form a prefix, and the position of the first
clear bit is therefore $\mathrm{popcount}(m)$. That value is exactly the lower
bound: the number of keys in the node strictly below the target. One compare, one
movemask and one population count give the answer with no branch and no
trailing-zero count.

Sortedness is what makes the mask a prefix, but the kernel never reads the mask's
shape, only its population count. Section~\ref{sec:append} uses that: the node the
shipped kernel actually reads is \emph{not} sorted in memory, and the same
population count is still the same answer.

Two details in the kernel are necessary for correctness and easy to get wrong. AVX2 has no
unsigned 32-bit compare, so \texttt{\_mm256\_cmpgt\_epi32} is used with both
operands biased by \texttt{0x80000000} to map unsigned order onto signed order.
Without the bias the empty-slot sentinel \texttt{0xFFFFFFFF} reads as $-1$ and
sorts below every real key, which inverts the padding that keeps a partially
filled node sorted. AVX-512 provides \texttt{\_mm512\_cmplt\_epu32\_mask}
natively and needs no bias.

Absence is reported as $k$, the node's key count, rather than as $-1$. The
distinction is not stylistic. A $-1$ sentinel requires a data-dependent
conditional to produce, and at a realistic hit ratio that conditional is
unpredictable. Measured across that change on the same harness at an earlier
revision, it cost 0.5 branch mispredictions per probe and roughly three quarters of
the AVX2 kernel's cycles, 25.4 against 6.2; both figures are in the artifact's
Phase 1 results rather than in Table~\ref{tab:kernels}, which reports the current
kernels only. Folding absence into the compare mask
removes the conditional, and for the ordered kernel the convention falls out of the
population count for free.

\subsection{Prefix-relative surrogates}
\label{sec:surrogate}

A vector lane holds 32 bits and a RocksDB key does not, so the lane holds a
surrogate. The surrogate must preserve order, which rules out a hash: a hash
distributes well and destroys the property the lower bound depends on. So the lane
holds four key bytes loaded big-endian, which makes an unsigned 32-bit compare of
two surrogates reproduce RocksDB's default bytewise \texttt{memcmp} order. Keys
shorter than four bytes pad with \texttt{0x00}, because a prefix sorts before any
key extending it.

\begin{table}[t]
\caption{Expected full-key comparisons per lookup, $\sum_i m_i^2 / N$ over
         surrogate multiplicities, at $N = 200{,}000$ keys. 1.0 is a perfect lane;
         $N$ itself means the lane holds one value for the whole keyspace.}
\label{tab:surrogate}
\centering\small\setlength{\tabcolsep}{4pt}
\begin{tabular}{lrr}
\toprule
key distribution & absolute & prefix-relative \\
\midrule
uniform random binary      &      1.0 & 1.0 \\
hex UUID                   &      4.0 & 1.0 \\
tenant-prefixed            &    201.0 & 1.1 \\
big-endian sequential      & 200000.0 & 1.0 \\
big-endian timestamp       & 200000.0 & 1.0 \\
decimal with a fixed stem  & 200000.0 & 1.0 \\
single table prefix        & 200000.0 & 1.0 \\
\texttt{db\_bench} default & 200000.0 & 1.0 \\
\bottomrule
\end{tabular}
\end{table}

Taking those four bytes from the start of the key does not work, and measurement
is what established it. Table~\ref{tab:surrogate} reports the expected number of
full-key comparisons a lookup performs, over eight key distributions at 200{,}000
keys each. An absolute surrogate takes a single value across the entire keyspace in
five of them, because a table name prefix, a tenant identifier or a big-endian
timestamp puts identical bytes at the front of every key in a database. A lane
holding one value for every key discriminates nothing, and the SIMD kernel
degenerates into a linear scan confirmed key by key. Only uniformly random binary
keys leave an absolute lane effective, and no one writes those.

The fifth of the five deserves naming, because it is this paper's own evaluation
keyspace. \texttt{db\_bench} writes the key ordinal big-endian into the first eight
bytes of a sixteen-byte key and pads the rest, so at the two million keys of
Section~\ref{sec:evaluation} the ordinal never reaches $2^{21}$ and the leading
five bytes are zero on every key in the database. Every RocksDB figure we report is
measured on a keyspace where an absolute surrogate would have held one value
throughout. That is the hardest case for the lane rather than the easiest, which is
worth stating plainly: the read results in Section~\ref{sec:evaluation} are not
obtainable without the correction described next.

The correction is that the lane never needed to be order-preserving globally, only
within a node. Each node stores the length of the prefix shared by all its keys,
and surrogates are taken immediately after it. Two consequences live in the code.
The prefix is recomputed from the first and last key whenever a node splits, since
an insert at either end can shorten it and a stale prefix mis-orders every
surrogate in the node. And a key that does not share the node's prefix has no
meaningful surrogate at that offset, so it is placed by full comparison instead;
such a key is provably outside the node's range rather than inside it, so its lane
saturates in the direction it sorts. The right-hand column of
Table~\ref{tab:surrogate} is that design, and it restores every distribution to an
effective lane, the worst at 1.1 comparisons per lookup.

The surrogate is a candidate and never an answer. Bytewise order on the user key
is not RocksDB's order in general: a column family may install a custom
comparator, and two entries whose user keys are equal are ordered by a
\emph{descending} sequence number packed into the internal key. Every SIMD
candidate is therefore confirmed against the comparator, with a binary search
fallback when confirmation fails. Four deliberately introduced kernel bugs passed
a full differential test against RocksDB's skip list because of this
confirmation, which is why the kernels also carry direct unit tests against a
scalar reference.

\subsection{Append-only nodes and the order word}
\label{sec:append}

An ordered node must stay sorted, and an earlier version of this design kept it
sorted by rebuilding it: every insert allocated a fresh payload, merged the new
key into the sorted run, and published the new payload pointer with one release
store. Publication was atomic and cheap to reason about. It also charged a whole
node to the arena for every key written, 445.4 bytes per key against a live
structure of 40, and scattered the nodes a lookup walks across ten times the
memory they occupy.

The current design stops moving keys. A slot is written once and never again, and
what gets republished is the \emph{order} over the slots. Sorted rank $i$ is held
in nibble $i$ of a 64-bit word, so a new order is one aligned 64-bit store and
publication stays exactly as atomic as swapping a pointer was. An insert is two
stores into a free slot, then one release store of the new order word. Nothing
already visible moves, so a reader that takes one acquire load of the order word
holds a consistent view: everything the word names is frozen.

A nibble holds $\mathit{slot}+1$ rather than the slot, so zero means ``this rank is
unused''. That costs one of the sixteen slots and buys two properties. The word
becomes self-describing, since the count is the position of its highest nonzero
nibble, so nothing needs a second store beside it; a second field would need a
second store and would reopen the torn states the single store exists to prevent.
And an unused rank decodes to slot 15, the lane no key occupies and every node
keeps at the sentinel, so a node is padded above its last key for free.

Slot order is insertion order, so lanes are not sorted in memory. The ordered
kernel turns out not to care. The lower bound is a count of the live keys below
the target, and a count over a set does not depend on how the set is arranged, so
permuting the lanes into rank order produces a different vector and the same
population count. An earlier version of this design did permute, at a measured
cost of 21.03 against 9.89 cycles per probe on AVX2, and removing the permutation
changed no answer.

What the kernel does still need is a mask over the live lanes, because a writer
may be storing into the slot at index \emph{count} while the kernel runs and an
unmasked 64-byte load would race with it. That mask is the entire residual cost of
the append-only layout, 1.50$\times$ on AVX-512 and 1.94$\times$ on AVX2 against
the same kernel over a sorted node. It is justified by the memory model \cite{boehm2008memorymodel} alone: no test
detects its removal, and ThreadSanitizer cannot either, because the compiler does
not instrument the vector load. The mask is a correctness device and not a
performance hint, because masked lanes in \texttt{vmaskmovps} and
\texttt{vmovdqu32} are architecturally not accessed at all; the unmasked load
would be a real data race that our tests happen not to observe.

\subsection{Descent hints}

Reaching the right node is a skip-list tower above level 0, with branching factor
four. Everything in the tower is an accelerator and never the truth, because the
search only follows a pointer to a node whose first key is at or below the target,
and nodes are ordered and never removed; a stale or missing tower pointer costs a
longer walk and never a wrong answer.

The descent was the dominant cost of a lookup and its size had been estimated
rather than measured. Measured, a lookup over 500{,}000 keys examined 30.0
successors, not the eight an earlier design note had assumed, and each one was a
chain of dependent loads ending in a virtual comparator call on a full internal
key. Each node now caches the first eight bytes of its first key in its header, so
a hop compares two integers in a line the hop has already loaded, and only a tie
falls through to the comparator.

Three invariants hold this up. The hint is immutable for every node but the head,
because the descent returns the last node whose first key is at or below the key
being inserted, so an insert into a non-head node sorts at position 1 or later and
a split leaves the left half's first key alone; the head is the exception and is
also the only node no hop ever moves to, so its hint is never read. A tie decides
nothing and falls through to the comparator, which is why the hint is eight bytes
and the surrogate four: a surrogate tie costs one comparison inside a node already
in cache, and a hint tie costs the miss chain the hint exists to avoid. And the
hint is stored before the payload, because a first key never rises, so a new hint
beside an old payload is a hint that is too small, and too small only ever lets a
hop enter a node it was already entitled to enter. The other order would expose a
hint that is too large, which would make a hop stop one node short of the key it
wanted.

The fast path is disabled unless the representation can confirm that the user
comparator agrees with bytewise order, because a wrong hop is not a slow lookup
but a lost key.

\subsection{Concurrency}

Readers are lock-free and never write. A reader takes one acquire load of the
order word and reads only slots that word names, all of which are frozen.

Writers take a per-node spin lock that pauses 64 times and then yields. This is a
deliberate narrowing of the lock-free goal \cite{herlihy2008art} and we state it
rather than hide it. A fully lock-free ordered insert has to shift keys inside a
sorted node, which is
not one atomic operation, so the design publishes through a single release store
with a retry and uses a lock to stop writers livelocking against each other.
Contention is per node, so two writers meet only when they target the same fifteen
keys. The backoff yields rather than spinning flat because RocksDB routinely runs
more writers than cores, where a flat spin burns every other thread's time slice
while the lock holder waits to be scheduled.

A split is the one operation that changes a node's contents and its successor at
the same time, so the level-0 link and the order word both live inside the
payload. The right half is published as a new node first, then a single release
store swaps the left half's payload pointer and its successor together. A reader
sees either the old full node or the new pair, never a state where a key is
duplicated across both or missing from both.

Section~\ref{sec:append} reduced what the writer lock covers from an allocation
and a fifteen-key merge to three stores, which makes replacing it with a
compare-and-swap on the order word a live question. We did not attempt it here,
for a reason Section~\ref{sec:tsan} makes concrete.

\section{Implementation}

Aparajita is a header-only C++20 library of roughly 2{,}000 lines plus a RocksDB
plugin. The core is a template over a traits policy carrying the three things that
differ between the standalone structure and the RocksDB representation: what an
entry is, how order is decided, and where memory comes from. The split, the
append path and the tower exist once and are shared by both.

The plugin derives from \texttt{MemTableRep} and \texttt{MemTableRepFactory} and
builds through RocksDB's \texttt{plugin/} mechanism, so \texttt{db\_bench
--memtablerep=aparajita} selects it without patching RocksDB sources. It
implements \texttt{Allocate}, \texttt{InsertKey} and its concurrent variants,
\texttt{Get}, \texttt{GetIterator}, \texttt{ApproximateMemoryUsage} and
\texttt{MarkReadOnly}. \texttt{IsInsertConcurrentlySupported} returns true, which
matters because RocksDB routes concurrent writes only to a representation that
advertises support; without it every multi-threaded write benchmark serialises,
which is part of why VectorRep loses on write throughput despite a friendlier
layout.

Two integration details are worth recording. \texttt{MemTableRep::Allocate} hands
key storage to the caller, so the representation never copies a key and an entry
is a bare 8-byte pointer into RocksDB's arena. And the descent hint's fast path
needs to confirm the comparator is the default bytewise one, for which
\texttt{MemTableRep::KeyComparator} exposes no route other than a
\texttt{dynamic\_cast}; RocksDB compiles Release with \texttt{-fno-rtti} by
default, so the evaluation builds with \texttt{USE\_RTTI=ON}, a supported CMake
option applied to the whole tree. The skip-list rows in
Section~\ref{sec:evaluation} are the control showing that flag alone moves
nothing.

Correctness is checked by a differential test against RocksDB's own skip list:
forward and backward iteration, every \texttt{Get}, \texttt{Seek} and
\texttt{SeekForPrev} agree byte for byte, under concurrent writes, across a flush,
and under a reverse comparator. RocksDB's own MemTable gtests build their options
through \texttt{CurrentOptions()} and never consult \texttt{--memtablerep}, and
\texttt{db\_stress} resolves that flag through a hardcoded enum, so neither can be
pointed at a plugin representation without patching RocksDB. That is a limitation
of the harness rather than of the plugin, and it is why a differential test stands
in for the upstream suite.

\section{Evaluation}
\label{sec:evaluation}

\subsection{Method}

All figures come from a GCP \texttt{c4-standard-24} in \texttt{us-central1-b}: a
Xeon Platinum 8581C (Emerald Rapids), 12 physical cores and 24 logical with SMT
on, 88\,GiB, Ubuntu 24.04, GCC 14.2, RocksDB v9.11.2 built Release with
\texttt{USE\_RTTI=ON}. One host, one RocksDB tree, one set of build flags, one
uninterrupted session, nothing else running. Every representation is selected
through \texttt{--memtablerep} on the same \texttt{db\_bench} binary, so the only
thing differing between columns is the representation.

The workload is \texttt{db\_bench}'s default 16-byte key and 100-byte value over a
2{,}000{,}000-key space, with compression off, automatic compaction off, the
write-ahead log disabled except where Section~\ref{sec:writes} enables it, and a
4\,GiB write buffer so that nothing flushes mid-run and the MemTable stays the
object under test. Every run reports its SST count, and no run behind
Tables~\ref{tab:reads} or \ref{tab:writes} produced more than the single file the
flush on close writes. Section~\ref{sec:flush} relaxes the write buffer to 64\,MiB
and is the one measurement here that flushes on purpose.

Five repetitions per cell. \texttt{fillrandom} runs in two independent passes,
because the scaling pass and the read passes each fill first, so a fill cell
carries ten samples and a read cell five.

Quota capped the shape at 24 vCPUs. The 1- and 4-thread points sit below physical
core count, 16 sits between physical and logical, and 64 is 2.7$\times$
oversubscribed. The top of every curve is contention rather than headroom and is
labelled as such throughout.

Two methodological points cost us a full run each and generalise past this
project.

\emph{A benchmark knob that sets both the work and the data has to be split before
a curve over it means anything.} \texttt{db\_bench}'s \texttt{--num} is per thread
and controls both how much work a thread does and the range its random keys are
drawn from. The natural way to hold total work constant across thread counts,
\texttt{--num=TOTAL/threads}, also divides the keyspace: at 1 thread the MemTable
held 1.26 million distinct keys and at 16 threads it held 125{,}000, each written
sixteen times, and the measured read hit rate drifted from 63\% to 100\% across
the curve. That is a working-set sweep in a thread-scaling curve's clothes, and it
flattered the reads. In the figures reported here \texttt{--num} is pinned at
2{,}000{,}000 and \texttt{--writes} and \texttt{--reads} carry the per-thread work,
holding the hit rate at 63\% everywhere.

\emph{Counter differencing has to cancel the setup exactly, not approximately.}
\texttt{perf} wraps a whole process, so a read benchmark's counters include the
fill that populated the MemTable, the database open and the flush on close.
Subtracting a separate fill-only run fails, because the fill is twenty times the
size of the read pass and the subtraction is two large noisy numbers differenced
to a small one; it produced negative instruction counts. Every counter figure
below is instead the difference between two runs of the \emph{same} benchmark that
differ only in operation count, with an identical fill at an identical seed, so
everything that is not the extra operations cancels exactly. Reads difference
20{,}000{,}000 against 2{,}000{,}000 and writes 2{,}000{,}000 against 200{,}000.

\subsection{Search kernels in isolation}

Table~\ref{tab:kernels} reports the kernels over 200{,}000 probes at a 0.5 hit
ratio, nine repetitions, median, with hits landing on a uniformly random slot and
the hit and miss mix randomised rather than blocked. Blocking the mix or always
probing slot 0 lets the branchy baseline exit after one comparison and hides
exactly the mispredictions the design removes.

\begin{table}[t]
\caption{Search kernels, cycles per probe, hit ratio 0.5}
\label{tab:kernels}
\centering\small\setlength{\tabcolsep}{4pt}
\begin{tabular}{lrrrr}
\toprule
kernel & cyc/probe & ins/probe & br/probe & miss/probe \\
\midrule
\multicolumn{5}{l}{\emph{equality search over a sorted node}} \\
scalar linear      & 27.65 & 34.50 & 15.24 & 0.5253 \\
scalar binary      & 70.24 & 51.41 & 13.16 & 1.9095 \\
scalar branchless  & 16.53 & 90.00 &  3.00 & 0.0000 \\
AVX2               &  5.01 & 20.00 &  3.00 & 0.0000 \\
AVX-512            &  4.04 & 15.00 &  3.00 & 0.0000 \\
\midrule
\multicolumn{5}{l}{\emph{lower bound over a sorted node}} \\
scalar             & 40.14 & 26.69 & 11.33 & 1.0390 \\
scalar branchless  &  8.67 & 44.00 &  3.00 & 0.0000 \\
AVX2               &  5.10 & 25.00 &  3.00 & 0.0000 \\
AVX-512            &  4.03 & 16.00 &  3.00 & 0.0000 \\
\midrule
\multicolumn{5}{l}{\emph{lower bound over an append-only node (shipped)}} \\
scalar             & 49.49 & 87.35 & 21.66 & 0.9276 \\
AVX2               &  9.89 & 41.00 &  4.00 & 0.0000 \\
AVX-512            &  6.04 & 30.00 &  4.00 & 0.0000 \\
\bottomrule
\end{tabular}
\end{table}

The branch behaviour is the cleanest result in the paper. Every vector kernel and
every branchless scalar kernel records 0.0000 mispredictions per probe, against
0.5253 for the branchy scalar baseline, and the vector kernels do not respond to
the hit ratio at all: at a 0.9 hit ratio every ordered row reproduces to within
0.02 cycles, while the branchy baseline moves from 28.49 to 40.47.

Two rows deserve a reader's scepticism and we state them before drawing any
conclusion from the table.

The shipped AVX2 ordered kernel, at 9.89 cycles per probe, is 14\% slower than
scalar branchless lower bound at 8.67. SIMD losing to scalar on the ISA that ships
would be a serious problem if the scalar kernel were an available alternative, and
it is not: it reads a sorted node, and no scalar branchless kernel exists for the
append-only layout. The comparison the design actually faces is 9.89 against
49.49. The wider point stands, which is that the append-only layout bought its
write win at a read-side price of 1.94$\times$ on AVX2 and 1.50$\times$ on
AVX-512, and that price is the live-lane mask described in
Section~\ref{sec:append}.

The equality rows for AVX2 and AVX-512 are not stable across our runs. A second
run on a different instance size of the same part reports 4.02 and 5.02 rather
than 5.01 and 4.04, inverting which ISA is faster while retained instruction
counts stay at 20 and 15. The ordered families reproduce across those same two
runs to within 0.2 cycles, so we quote the ordered rows and make no claim about
the relative speed of the two ISAs on equality search.

\subsection{Point lookups}

Table~\ref{tab:reads} reports \texttt{readrandom} over a MemTable-resident
working set with the write-ahead log disabled, and \texttt{readwhilewriting} with
one writer alongside the readers. No run in either table produced more than the
single SST file that the flush on close writes, so every figure is served by the
MemTable.

\begin{table}[t]
\caption{Read throughput, kops/s, WAL off. Median of five runs; the rng column is
         $(\max-\min)/\mathrm{median}$ over those five. VectorRep is four orders of
         magnitude slower and is reported in Section~\ref{sec:vectorrep}.}
\label{tab:reads}
\centering\footnotesize\setlength{\tabcolsep}{3pt}
\begin{tabular}{lrrrrr}
\toprule
threads & Aparajita & rng & skip list & rng & vs.\ skip list \\
\midrule
\multicolumn{6}{l}{\texttt{readrandom}} \\
1  &   731.9 &  7.8\% &  528.2 & 17.7\% & +38.6\% \\
4  &  2703.1 &  3.0\% & 1969.1 &  5.2\% & +37.3\% \\
16 &  8773.9 &  8.4\% & 7181.8 & 23.8\% & +22.2\% \\
64 & 10913.5 &  4.7\% & 8677.9 &  9.8\% & +25.8\% \\
\midrule
\multicolumn{6}{l}{\texttt{readwhilewriting}} \\
1  &   705.6 &  2.6\% &  519.7 &  4.4\% & +35.8\% \\
4  &  2693.8 &  2.7\% & 1950.2 &  6.9\% & +38.1\% \\
16 &  8691.9 & 31.8\% & 6738.8 &  8.4\% & +29.0\% \\
64 & 10531.9 &  7.8\% & 8578.5 & 13.3\% & +22.8\% \\
\bottomrule
\end{tabular}
\end{table}

Aparajita is ahead at every thread count on both benchmarks, and at seven of the
eight points the samples do not overlap at all: Aparajita's slowest of five runs
beats the skip list's fastest. The exception is \texttt{readwhilewriting} at 16
threads, where one Aparajita run came in at 6{,}082\,kops/s against a skip-list
best of 7{,}019, which is also the widest range in the table at 31.8\%. The other
seven cells separate cleanly, and the two benchmarks agree with each other at every
thread count.

Table~\ref{tab:counters} gives the marginal per-operation cost from the two-point
differencing described above. Aparajita retires 3{,}271.5 instructions per lookup
against the skip list's 7{,}349.9, and takes 62.18 L1 misses against 103.99. Those
are the two mechanisms the design targets, showing up where they were predicted:
fewer instructions from replacing a chain of virtual comparator calls with integer
compares in already-loaded lines, and fewer L1 misses from a structure that
consumes 100.6 rather than 445.4 bytes of arena per key.

\begin{table}[t]
\caption{Marginal cost per operation, 16 threads}
\label{tab:counters}
\centering\footnotesize\setlength{\tabcolsep}{2.5pt}
\begin{tabular}{llrrrr}
\toprule
benchmark & rep & instr & br miss & br miss \% & L1 miss \\
\midrule
\texttt{fillrandom} & Aparajita & 22181.7 & 16.15 & 0.361 & 178.86 \\
                    & skip list & 24177.1 & 18.01 & 0.379 & 227.73 \\
                    & VectorRep & 14989.5 & 50.70 & 1.636 & 176.96 \\
\midrule
\texttt{readrandom} & Aparajita &  3271.5 &  8.46 & 1.241 &  62.18 \\
                    & skip list &  7349.9 & 10.39 & 0.727 & 103.99 \\
\midrule
\texttt{readwhilewriting} & Aparajita &  3622.2 &  9.53 & 1.265 &  72.41 \\
                          & skip list &  7597.6 & 11.93 & 0.808 & 110.70 \\
\bottomrule
\end{tabular}
\end{table}

One inversion in that table is worth stating rather than burying. Aparajita's
branch miss \emph{rate} on reads is worse than the skip list's, 1.241\% against
0.727\%, while its branch misses \emph{per operation} are better, 8.46 against
10.39. Both are true, and they are consistent: the design executes less than half
as many branches per lookup, so the same absolute count of mispredictions is a
larger share of a much smaller denominator. Per-operation cost is what a workload
pays. A paper quoting the rate alone would be quoting the single framing in which
this design looks worse than the structure it replaces.

\subsection{Ordered access}
\label{sec:seeks}

Section~\ref{sec:reps} argues that the design has to be ordered because a MemTable
iterator has to be, and Section~\ref{sec:vectorrep} measures what giving that up
costs VectorRep. Table~\ref{tab:seeks} measures what the ordering buys
and what it costs, and the answer divides at the boundary between finding a
position and walking from it.

\begin{table}[t]
\caption{\texttt{seekrandom} throughput, kops/s, WAL off, over a resident MemTable.
         Median of five runs. \emph{nexts} is the number of \texttt{Next} calls
         issued after each \texttt{Seek}.}
\label{tab:seeks}
\centering\footnotesize\setlength{\tabcolsep}{4pt}
\begin{tabular}{lrrr}
\toprule
threads & Aparajita & skip list & vs.\ skip list \\
\midrule
\multicolumn{4}{l}{\texttt{Seek} alone, nexts $= 0$} \\
1  &  428.2 &  333.0 & +28.6\% \\
4  & 1744.3 & 1368.8 & +27.4\% \\
16 & 5688.7 & 4853.0 & +17.2\% \\
64 & 7639.7 & 6636.1 & +15.1\% \\
\midrule
\multicolumn{4}{l}{\texttt{Seek} then ten \texttt{Next}, nexts $= 10$} \\
1  &  191.7 &  197.6 & $-$3.0\% \\
4  &  744.1 &  769.6 & $-$3.3\% \\
16 & 2246.6 & 2348.0 & $-$4.3\% \\
64 & 3093.8 & 3217.5 & $-$3.8\% \\
\bottomrule
\end{tabular}
\end{table}

A Seek on its own is 15\% to 29\% faster than the skip list at every thread count,
narrowing as threads rise exactly as the point lookups do, and for the same reason:
a Seek is a descent followed by a lower bound, which is what the hints and the
kernel make cheap.

A Seek followed by ten Nexts is 3\% to 4\% \emph{slower}. The two results do not
conflict. A skip-list \texttt{Next} is one pointer load along the level-0 chain.
Aparajita's decodes the next rank from the order word, indexes the slot that nibble
names, and follows a key pointer into the arena, which is more work per step for
precisely the reason the append-only layout is cheap to write: the sorted order is
a computed permutation and not the physical one. At ten steps per Seek the walk
dominates the descent and the design is behind. The margin is small and the ranges
overlap at 1, 4 and 64 threads; at 16 they do not, and the sign is the same at all
four points, so we report a loss rather than a tie.

That bounds the claim this paper makes. Aparajita is a faster place to look a key
up and a slightly slower place to scan from, and a workload dominated by long range
scans over the MemTable is not one these numbers support. The fix has an obvious
shape and we did not attempt it: the order word already sits in a line the iterator
has loaded, so a \texttt{Next} that decoded several ranks at once would amortize
what is currently paid per step.

\subsection{Inserts, and the ceiling \texttt{db\_bench} imposes}
\label{sec:writes}

Table~\ref{tab:writes} reports \texttt{fillrandom}. Aparajita is 18.8\% ahead at
one thread, 2.6\% ahead at four, and tied at 16 and 64. The 4-thread figure
reproduces an independent 2.0\% measured on a different host in an earlier phase.
The ties are in the samples and not a rounding of the medians. At 16 threads the
two independent fill passes disagree about the sign, giving $-3.0\%$ and $+0.0\%$,
and the ranges overlap outright: 1387.9--1430.1\,kops/s for Aparajita against
1397.2--1450.1 for the skip list. At 64 threads both lie inside a 1\% band,
503.2--507.8 against 500.2--504.3.

\begin{table}[t]
\caption{\texttt{fillrandom} throughput, kops/s, WAL off}
\label{tab:writes}
\centering
\begin{tabular}{lrrrr}
\toprule
threads & Aparajita & skip list & VectorRep & vs.\ skip list \\
\midrule
1  &  678.9 &  571.2 & 2222.9 & +18.8\% \\
4  &  677.9 &  660.9 & 1201.0 &  +2.6\% \\
16 & 1394.8 & 1437.7 &  911.1 &  $-$3.0\% \\
64 &  505.8 &  502.6 & 1101.0 &  +0.6\% \\
\midrule
\multicolumn{5}{l}{\emph{scaling relative to 1 thread}} \\
4  & 1.00$\times$ & 1.16$\times$ & 0.54$\times$ & \\
16 & 2.05$\times$ & 2.52$\times$ & 0.41$\times$ & \\
64 & 0.75$\times$ & 0.88$\times$ & 0.50$\times$ & \\
\bottomrule
\end{tabular}
\end{table}

The scaling rows are the tell. Neither representation scales, and two structures
with entirely different insert paths do not accidentally share a scaling curve.
The counters name the cause: a \texttt{fillrandom} operation at 16 threads retires
22{,}181.7 instructions in Aparajita and 24{,}177.1 in the skip list. A MemTable
insert is not twenty-two thousand instructions; an append into a node under a spin
lock is a few hundred. The rest is RocksDB's write-thread group, where writers
spin waiting for a group leader to commit the batch, charged identically to every
representation. At 16 threads on 12 cores, \texttt{db\_bench fillrandom} measures
RocksDB's write path with the MemTable as a rounding error.

\texttt{memtablerep\_bench} has no write group and no SST, block cache or version
machinery, and the difference reappears intact. Insert costs 0.5435\,\textmu s/op
against the skip list's 0.8040, the skip list 48\% slower, and the gap holds at
48\% to 52\% across all four
thread counts. RocksDB hardcodes that benchmark's \texttt{fillrandom} to a single
thread regardless of \texttt{--num\_threads}, so those four columns are four
samples of one single-threaded measurement rather than a curve; they agree to
within 3\%, which is the useful thing they say. Its \texttt{readrandom} narrows under oversubscription, from 62.5\% at one thread
to 10.4\% at 16 and 0.4\% at 64, which looks like a contradiction of
Table~\ref{tab:reads}'s +25.8\% at 64 threads in the harness with more machinery in
it. The absolute figures reconcile the two. At 64 threads
\texttt{memtablerep\_bench} costs 52.2 and 52.5\,\textmu s per lookup for the two
representations, against \texttt{db\_bench}'s 4.4 and 5.7: the rep-level harness has
degraded a hundredfold from its own single-threaded 0.50\,\textmu s while
\texttt{db\_bench} degrades threefold. A lookup costing on the order of a
microsecond is a fiftieth of a 52\,\textmu s budget, so at 64 threads that benchmark
has stopped resolving the representation rather than the representations having
converged. We read its 1- and 4-thread points and not its 64-thread one, and the
read claim in this paper rests on Table~\ref{tab:reads}.

The claim this evidence supports is that the representation inserts roughly half
again as fast as the skip list, visible undiluted at one thread and in
\texttt{memtablerep\_bench}, and swamped by RocksDB's write-group machinery in
multi-threaded \texttt{db\_bench}. It does not support a multi-threaded
\texttt{db\_bench} write scaling claim, and we do not make one.

The write-ahead log behaves as expected and we measured it rather than assuming
it. At 16 threads with the log enabled the two representations run at 1043.0 and
1048.1\,kops/s against 1394.8 and 1437.7 with it disabled: the log costs about
27\% and closes the already-small gap to 0.5\%.

\subsection{Arena cost, and what it costs to flush}
\label{sec:flush}

Every figure above holds the MemTable resident, which is what isolates the
representation and is also what hides its one clear regression. Aparajita charges
about 1.4 times the skip list's arena per key. Table~\ref{tab:flush} makes that
visible by shrinking the write buffer to 64\,MiB and letting the same two million
keys flush: Aparajita fits 285{,}714 keys into a MemTable against the skip list's
400{,}000 and writes seven L0 files where the skip list writes five.

\begin{table}[t]
\caption{Arena cost made visible: 64\,MiB write buffer, one thread, 2{,}000{,}000
         keys. A representation charging more arena per key fits fewer keys into a
         MemTable and flushes more often.}
\label{tab:flush}
\centering\footnotesize\setlength{\tabcolsep}{4pt}
\begin{tabular}{lrrrr}
\toprule
rep & fill \textmu s/op & L0 files & keys/flush & flush \textmu s/key \\
\midrule
Aparajita & 1.031 & 7 & 285{,}714 & 0.425 \\
skip list & 1.313 & 5 & 400{,}000 & 0.441 \\
\bottomrule
\end{tabular}
\end{table}

Flushing is not what those extra files cost. The same two million keys are written
either way and the per-key flush cost is within four percent, 0.425 against
0.441\,\textmu s, so the total time spent flushing is if anything slightly lower.
Walking the structure, which is what a flush does, is not where the arena overhead
lands, and the fill itself is 21\% faster here because this configuration is the
one where the write buffer is small enough for the insert path to matter.

What the overhead costs is L0 files, forty percent more of them for the same data.
That is compaction work, which this configuration disables and therefore does not
measure, and it is also more frequent write stalls at any realistic
\texttt{level0\_slowdown\_writes\_trigger}. We state it as an open cost rather than
estimate it.

\subsection{VectorRep}
\label{sec:vectorrep}

VectorRep has the friendliest memory layout of the three and loses anyway, which
is the comparison that motivates an ordered design. It is by far the fastest
single-threaded fill at 2{,}222.9\,kops/s, because appending to a vector is the
cheapest insert available, and it is the only representation whose fill throughput
\emph{falls} as threads are added, 0.54$\times$ at 4 and 0.41$\times$ at 16,
because \texttt{IsInsertConcurrentlySupported} returns false and every write
serialises.

Its reads are not slow so much as disqualifying: 0.9 to 3.3 seconds per lookup,
and $6.4 \times 10^9$ retired instructions per lookup.
\texttt{VectorRep::Get} builds an iterator, and on a mutable MemTable that
iterator copies the entire bucket and sorts the copy, because a sorted order
cannot be cached on a structure still being appended to. Its read benchmarks are
capped at 20 operations for this reason and reported as capped; uncapped they
would have run for days.

\subsection{Race freedom}
\label{sec:tsan}

The concurrent insert and read test runs under ThreadSanitizer
\cite{serebryany2009tsan} at 1, 4, 16 and 64 threads, and reports no races at any
thread count on any run.

The cost of that instrumented run is heavy-tailed in a way worth recording,
because we published a wrong figure for it before measuring it properly. Five runs
of the same binary on one idle host gave totals of 9.4\,s, 186\,s, 298\,s, 373\,s
and 524\,s for identical work, and a single 6.9\,ms sample from the fast tail had
been reported as the result. The per-phase breakdown locates it: 640 total inserts
take under 3\,ms of race time at 4 threads and up to nine minutes at 64. What
grows is the number of threads spinning in an instrumented \texttt{test\_and\_set}
loop, where ThreadSanitizer's per-thread vector clocks make each of those atomics
cost $O(\text{threads})$.

That is an artifact of instrumentation and not evidence about uninstrumented
throughput, and we are explicit that it is not a reason to believe the spin lock
is expensive in production. It is a reason to want the compare-and-swap
formulation, since that removes the loop the instrumentation chokes on.

\section{Limitations}

Three exit criteria this work set itself are not fully met, and one of them
concerns the mechanism the design is most directly about.

\textbf{Cross-core invalidation is unmeasured.} Loads that hit a line another core
holds modified are the counter that would speak most directly to a cache-line
argument, and we could not collect it. The cloud provider's API rejected the vPMU
tier documented to cover LLC events on both API versions tried. Under the tier it
did accept, \texttt{mem\_load\_l3\_hit\_retired.xsnp\_fwd} enumerates in
\texttt{perf list}, programs without error, reports 100\% enabled time, and then
counts exactly zero against a run retiring $1.08 \times 10^{11}$ instructions
across 16 threads. \texttt{cache-references} and \texttt{cache-misses} return
\emph{not supported} outright. These are PEBS-backed events and the hypervisor
does not implement PEBS \cite{intel-sdm}, so the counter reads zero rather than
refusing.

That failure mode is worse than an unsupported event, because ``0.000 cross-core
hits per operation'' is a publishable-looking number and it would have been false.
Our harness now calibrates the event against a workload that must produce
cross-core sharing and demotes it to unavailable if it stays at zero. Two related
traps are worth passing on: Intel retired the \texttt{XSNP\_HITM} name at Ice Lake
and the counterpart on this part is \texttt{XSNP\_FWD}, so resolving the event
name against the host's own \texttt{perf list} matters; and \texttt{perf list |
grep -q} makes \texttt{grep} exit at its first match, \texttt{perf list} then dies
of \texttt{SIGPIPE}, and under \texttt{set -o pipefail} every event the host has is
reported absent. Closing this needs bare metal or a provider that exposes PEBS to
guests.

\textbf{The scaling curve has no headroom.} Quota capped the host at 12 physical
cores, so every point above 12 threads is contention. That regime is arguably the
one RocksDB runs in, since more writers than cores is normal, which is why the
per-node backoff yields rather than spinning flat. It still cannot be read as
scaling, and the shape of the curve on a machine with 64 real cores is unknown.

\textbf{The write argument depends on a benchmark most readers do not use.}
Because multi-threaded \texttt{db\_bench fillrandom} measures RocksDB's write path
rather than the representation, the insert win has to be argued from
\texttt{memtablerep\_bench} and from the single-threaded point.

\textbf{Two cache effects adjacent to the argument are unmeasured.} The design
aligns to 64 bytes and pads \texttt{ListNode} to a full line so that writers to
different nodes do not share one. Intel parts also run an adjacent-line prefetcher
that pulls lines in pairs, so two nodes landing on lines $N$ and $N{+}1$ can still
be invalidated together under concurrent writers, and we did not isolate that: it
would need the same cross-core counter that came back dead, which is why it is a
gap rather than an oversight. Separately, the host is a single socket, so nothing
here says how the structure behaves when a writer's arena block is resident on
another NUMA node. Neither effect is under the representation's control in RocksDB,
where node memory comes from \texttt{ConcurrentArena}, but both would matter on a
two-socket machine.

\textbf{The consequences of flushing more often go unmeasured.} The arena
overhead in Section~\ref{sec:flush} produces forty percent more L0 files for the
same data. Flush time absorbs none of that, but compaction and the write stalls
that L0 file counts trigger would, and every configuration here disables automatic
compaction so that the MemTable is what is being measured. An end-to-end comparison
under compaction is the natural next measurement and would be the first place this
design could lose overall while winning on every number reported above.

Two smaller omissions are deliberate. The representation does not implement the
prefix-extractor and Bloom filter paths \cite{bloom1970,dayan2017monkey}, so it
inherits the base class's full iterator, which is correct but gives up the skip
list's prefix-Bloom short circuit. And AVX-512 is opportunistic behind a CPUID
check rather than required, on availability grounds rather than on the frequency
behaviour the optimization manual describes for wide vector work
\cite{intel-optimization}: AVX-512 is absent from most AMD parts before Zen 4, and
Intel disabled it in client parts after Rocket Lake.

\section{Conclusion}

Replacing a pointer-chasing skip list with cache-line nodes searched by a
branchless SIMD kernel improves MemTable point lookups by 22\% to 39\% inside
RocksDB, with the improvement traceable to 55\% fewer retired instructions and
40\% fewer L1 misses per lookup. The improvement is in finding a key and not in
walking from one: an ordered seek is 15\% to 29\% faster and a seek followed by ten
iterator steps is 3\% to 4\% slower, because the sorted order over an append-only
node is a permutation to be decoded rather than a chain to be followed. The same
trade shows up in space, where the representation charges 1.4 times the skip list's
arena per key. The kernel work is necessary and far from
sufficient: two structural problems outside the kernel, an expensive tower descent
and a per-insert node rebuild, each cost more than the kernel saved, and the
design only became competitive once both were fixed.

The result we did not get is as informative as the ones we did. Multi-threaded
\texttt{db\_bench} inserts are bounded by RocksDB's write-thread group at over
22{,}000 instructions per operation, which no MemTable design can move, and the
counter that speaks most directly to cross-core invalidation could not be
collected on virtualised hardware at all. Both are properties of the measurement
environment rather than of the structure, and both are the kind of thing that is
easier to state than to discover.

\section*{Availability}

Source, measurement scripts and the raw output behind every figure are at
\url{https://github.com/sinhaparth5/aparajita-memtable}, under Apache 2.0 or MIT
at the user's option \cite{aparajita}.

\bibliographystyle{IEEEtran}
\bibliography{refs}

\end{document}
